
\documentclass[12pt,a4paper]{article}
\usepackage[brazilian]{babel}
\usepackage{amsmath,amssymb,amsthm}
\usepackage{geometry}
\usepackage{enumitem}
\usepackage{hyperref}
\usepackage{xcolor}
\usepackage{graphicx}
\usepackage{fancyhdr}
\usepackage{tcolorbox}
\usepackage{totpages}
\usepackage{titlesec}
\usepackage{tasks}

% Configuração global das alternativas com o pacote tasks
\settasks{
  label=\textbf{\Alph*)},
  label-width=1.5em,
  label-offset=0.3em,
  item-indent=0pt,
  column-sep=1.2em
}

% Margens ajustadas para cabeçalho institucional
\geometry{top=3.5cm, bottom=2.5cm, left=2.0cm, right=2.0cm, headheight=2.6cm}

% Paleta de Cores
\definecolor{primary}{RGB}{10, 45, 90}       % Azul Escuro Institucional
\definecolor{secondary}{RGB}{25, 110, 180}   % Azul Destaque
\definecolor{bglight}{RGB}{245, 247, 250}    % Cinza Claro para Caixas
\definecolor{borderlight}{RGB}{210, 220, 230}% Borda Suave
\definecolor{correctgreen}{RGB}{40, 140, 60}  % Verde para Gabarito

% Configuração das Caixas (tcolorbox)
\tcbuselibrary{skins, breakable}

% Configuração de Cabeçalho e Rodapé
\pagestyle{fancy}
\fancyhf{}

% Cabeçalho com Logo e Disciplina
\fancyhead[L]{%
  \includegraphics[height=2cm]{Imagens/logocaruaru.png}%
}
\fancyhead[R]{%
  \small\sffamily
  \textbf{Instituto Federal de Pernambuco -- Campus Caruaru} \\
  \textcolor{secondary}{Disciplina:} Matemática VI \\
  \textcolor{secondary}{Prof.:} Cleibson José da Silva
}

% Rodapé
\fancyfoot[L]{\small\sffamily\color{gray} 7 -- Equações Algébricas 2}
\fancyfoot[R]{\small\sffamily\color{gray}Página \thepage\ de \ref{TotPages}}

% Linhas divisórias
\renewcommand{\headrulewidth}{0.8pt}
\renewcommand{\headrule}{{\color{primary}\hrule width\headwidth height \headrulewidth}}
\renewcommand{\footrulewidth}{0.4pt}

\hypersetup{
    colorlinks=true,
    linkcolor=secondary,
    urlcolor=secondary,
    citecolor=primary
}

\begin{document}

% Cabeçalho de Identificação do Estudante
\begin{tcolorbox}[colback=bglight, colframe=primary, boxrule=1pt, arc=3pt]
  \sffamily\small
  \begin{tabular}{@{}p{0.65\textwidth} p{0.3\textwidth}@{}}
    \textbf{Estudante:} \hrulefill & \textbf{Data:} \hrulefill \\[0.25cm]
    \textbf{Curso/Turma:} \hrulefill & \textbf{Nota/Visto:} \hrulefill
  \end{tabular}
\end{tcolorbox}

% Título Principal da 
\begin{center}
  {\LARGE\bfseries\color{primary}  7: Equações Algébricas 2} \\[0.1cm]
  {\small\sffamily\color{gray} Teorema das raízes racionais e relações de Girard}
\end{center}

\vspace{0.1cm}

\begin{enumerate}
    \item \textbf{(FUVEST-SP)} Seja $p(x) = x^4 + bx^3 + cx^2 + dx + e$ um polinômio com coeficientes inteiros. Sabe-se que as quatro raízes de $p(x)$ são inteiras e que três delas são pares e uma é ímpar. Quantos coeficientes pares tem o polinômio $p(x)$?
    \begin{enumerate}
        \item $0$
        \item $1$
        \item $2$
        \item $3$
        \item $4$
    \end{enumerate}

    \vspace{0.3cm}

    \item \textbf{(UFJF-MG)} Seja $S$ a soma das raízes do polinômio $p(x) = ax^2 + bx + c$, em que $a$, $b$ e $c$ são números reais e $a \ne 0$. Se $S_1$ é a soma das raízes de $p(x - 1)$, então a diferença $S_1 - S$ é
    \begin{enumerate}
        \item $-1$
        \item $0$
        \item $1$
        \item $2$
    \end{enumerate}

    \vspace{0.3cm}

    \item \textbf{(FUVEST-SP)} Se a equação $8x^3 + kx^2 - 18x + 9 = 0$ tem raízes reais $a$ e $-a$, então o valor de $k$ é
    \begin{enumerate}
        \item $\dfrac{9}{4}$
        \item $2$
        \item $\dfrac{9}{8}$
        \item $-2$
        \item $-4$
    \end{enumerate}

    \vspace{0.3cm}

    \item \textbf{(UFSCar-SP)} Sabendo-se que a soma de duas das raízes da equação $x^3 - 7x^2 + 14x - 8 = 0$ é igual a $5$, pode-se afirmar a respeito das raízes que
    \begin{enumerate}
        \item são todas iguais e não nulas.
        \item somente uma raiz é nula.
        \item as raízes constituem uma progressão geométrica.
        \item as raízes constituem uma progressão aritmética.
        \item nenhuma raiz é real.
    \end{enumerate}

    \vspace{0.3cm}

    \item \textbf{(UFMG)} A soma de todas as raízes da equação $(x - 1)^2 - (x - 1)(x + 4) = (x - 1)(x + 1)$ é
    \begin{enumerate}
        \item $-5$
        \item $-2$
        \item $2$
        \item $5$
        \item $6$
    \end{enumerate}

    \vspace{0.3cm}

    \item \textbf{(UFMG)} Se a equação $x^2 + px + q = 0$ admite raízes reais simétricas, então
    \begin{enumerate}
        \item $p = 1$ e $q = 0$
        \item $p = 1$ e $q > 0$
        \item $p = 1$ e $q < 0$
        \item $p = 0$ e $q > 0$
        \item $p = 0$ e $q < 0$
    \end{enumerate}

    \vspace{0.3cm}

    \item \textbf{(FUVEST-SP)} Sabe-se que o produto de duas raízes da equação algébrica $2x^3 - x^2 + kx + 4 = 0$ é igual a $1$. Então, o valor de $k$ é
    \begin{enumerate}
        \item $-8$
        \item $-4$
        \item $0$
        \item $4$
        \item $8$
    \end{enumerate}

    \vspace{0.3cm}

    \item \textbf{(UFRGS)} Se os números $-3$, $a$ e $b$ são raízes da equação $x^3 + 5x^2 - 2x - 24 = 0$, então o valor de $a + b$ é
    \begin{enumerate}
        \item $-6$
        \item $-2$
        \item $-1$
        \item $2$
        \item $6$
    \end{enumerate}
\end{enumerate}

\textbf{Gabarito:} 1) D\quad 2) D\quad 3) E\quad 4) C\quad 5) A\quad 6) E\quad 7) A\quad 8) B\quad 
\end{document}